\documentclass[12pt]{article}
\usepackage{amsmath,amssymb}

\begin{document}
\section*{{2025 AMC 12A Problems/Problem 23}}
Source: AoPS Wiki

\subsection*{Solution}
To satisfy the conditions, a $\textit{fair}$ integer must have no digit be a local minimum. Let's say we have $n$ distinct digits, with each digit being a number from $1$ to $9$ . To create a $\textit{fair}$ integer, we begin by placing the largest digit. For the second-largest digit, we can either place this digit to the right or to the left of the string already created. We have these $2$ options for the third-largest digit, and so on. Therefore, there are $2^{n-1}$ valid permutations to create a $\textit{fair}$ integer. We must also choose which digits will be in the permutation. If you are creating an $n$ -digit long $\textit{fair}$ integer, there are $9\choose{n}$ ways to pick which digits will be in the number. Therefore, for each $n \in \{1,2,\dots, 9\}$ , the number of fair integers of length $n$ is: \[\binom{9}{n} \cdot 2^{n-1}.\] Summing over all $n$ : \[\sum_{n=1}^9{\binom{9}{n} \cdot 2^{n-1}}=\frac{1}{2}\left(\sum_{n=0}^9{\binom{9}{n}}2^n -1\right)=\frac{1}{2}\left((1+2)^9 -1 \right) = \frac{1}{2}(19682) = \boxed{9841}.\] Note that the Binomial Theorem was used to equate \[\sum_{n=0}^9{\binom{9}{n}}2^n = (1+2)^9.\] ~lprado

\end{document}
